\section{Extending SPRUCE to anonymous credentials}
\label{sec:anonymous-credentials}

We target ARC as an anonymous credential schemes for two reasons:
\begin{description}
    \item[Simplicity] ARC uses re-randomization to generate several tokens. In contrast, balance-based schemes like Anonymous Credit Tokens~(ACT)~\cite{schlesinger-cfrg-act-00} require homomorphic encryption to maintain encrypted balances. 
    \item[Parallel Spends] ARC presentations can be generated and spent concurrently. Each ARC presentation is independently verifiable without cryptographic state updates. Balance-based schemes like ACT require sequential spending to update the state inbetween spends. 
\end{description}

Our goal is a post-quantum variant of ARC that preserves these properties while providing lattice-based security.

\subsection{ARC Credential Structure}

The credential binds $m_1$, a random secret generated by the client, and $m_2$, a deterministic secret derived from the application context, to the credential. $m_1$ is used for the tag generation and binds all presentations to the same underlying credential. 

$m_2$ is a deterministic secret derived from the application context. Its concrete usage depends on the application, but it is a public value. It can bind the credential to specific application data (like a place where you can spend your credential). The seperation between $m_1$ and $m_2$ allows context-specific credentials without new issuance, 

\subsection{Credential issuance}

\subsubsection{Client Credential Request}
In classical ARC, the client samples two random scalars $r_1,r_2\xleftarrow{\$} \mathbb{Z}_p$  and creates Pedersen commitments to hide $m_1$ and $m_2$ during issuance. 

The post-quantum variant of this issuance using SPRUCE would be to sample $(m_1, m_2, r_1, r)$. The user uses $r$ to compute the Ajtai commitment $c=A\left(\begin{array}{c}m_1\\m_2\\r_1//r\end{array}\right)$ and a proof that: 
\begin{enumerate}
    \item $c$ is well-formed
    \item $m_1, m_2, r_1, r$ are all small lattice elements
\end{enumerate}

\lena{For me, the critical difference here is that we have one commitment instead of two used in the original ARC, where we have one to $m_1$ and one to $m_2$. I did think this was a bit weird in the original ARC, but it does have the advantage of context-specific credentials without new issuance. I think if we don't do this we may lose some functionality. 
}

\subsubsection{Server Credential Response}
To issue a credential in ARC, the server samples a randomization scalar $b \xleftarrow{\$} \mathbb{Z}_p$ and computes the randomized $U = b \cdot G$, which serves as the base for all other credential elements. It will be re-randomized by the client to get an unlinkable credential. Since $b$ is unknown to the client, it prevents credential forgery. 

The server computes additional blinding values and issues a blinded MAC which simplifies to $U' = b \cdot (x_0 + x_1 m_1 + x_2 m_2) \cdot G $ after the client removes blinding terms.
The final credential is $(m_1, U, U', X_1)$, where $X_1=x_1*H$, where $x_1$ is one of the server's private keys and $H$ is a generator. From this one credential, the client can generate up to $n$ unlinkable presentations.

In the lattice setting, we propose to add a challenge $r_2$ that the server sends to the user. This randomness $r_2$ has enough min\_entropy to hide the GPV trapdoor. The user sends $t=h_{m_1\parallel m_2,r_1+r_2}(B)$, which is signed by the server after the server verifies the consistency of the earlier messages. 
This adds another two rounds to the issuance, which is fine since the issuance does not happen too often. 
\lena{I'm not sure where $B$ is coming from, and how we can prove efficiently that the values are the same as in the commitment. Has this challenge approach been used in lattices before? It seems quite straightforward in any case. }

\subsection{Credential Presentation}
Each presentation re-randomizes the credential elements with fresh randomness and computes a tag $(m_1 + \text{nonce})^{-1} \cdot \mathsf{H}(\text{context}, \text{"Tag"})$. The presentation includes a zero-knowledge proof of: 
    \begin{enumerate}
        \item a valid signature possession
        \item a correct tag construction
        \item a nonce within range $[0, n)$
    \end{enumerate}
    Since different nonces produce different tags, the same credential with two different nonces creates two unlinkable presentations. 

    In the lattice setting, $F$ is a PRF. The user chooses a new nonce $n$ and sends $f=F(m_2,n)$ alongside a proof of a hash of a valid GPV signature that contains $m_2$, and is consistent with $n$. 
    For the post-quantum construction we of course include a proof that the user has knowledge of a valid preimage of $t$, which also proves $m_1$. 
    Only proving signature, nonce and valid signature and proof of nonce in range is not enough, since the proof of the signature is randomized. This means that for the same nonce we can generate several proofs, which is why we evaluate $F$. Keeping $m_2$ secret means that we cannot correlate the values. 
    \lena{jetlag, TODO reread later}

\subsection{Limiting the number of presentations}
The presentation limit in ARC is \emph{not} cryptographically enforced by the protocol, but by an application parameter. There are two options to enforce the presentation limit:

\begin{description}
    \item[Random Nonce Selection] 
    The client randomly selects unused nonces from $[0, n)$ and reveals each nonce alongside its presentation. The server then validates that the nonce is in $[0,n)$ and stores observed tags to prevent double-spending. The presentations cannot be linked unless the nonces are reused. However, the client has to keep track of the spent nonces to prevent accidental double-spending. 
    \item[Range proofs] Instead of revealing the nonce, the client adds a zero-knowledge range proof proing that the nonce is in $[0, n)$. using range proofs simplifies the client state, as the client does not have to keep track of spent nonces. 
    However, the additional zero-knowledge proof increases the size of the presentation and verification cost.
\end{description}
It is explicitly the server's responsibility to enforce its rate limit, meaning the server has to keep track of the nonces and proofs to prevent double-spending. 

\lena{
To construct an ARC-style protocol from a blind signature scheme, we need to add the proofs mentioned in the presentation section and figure out a way to prove both the nonce and the signature. It would be interesting to see which approach is easier (range proofs or random nonces). }

\subsection{Considerations}
\begin{description}
    \item[Verifiability] ARC and other schemes currently uses privately-verifiable MACs (i.e. the issuer must be the verifier, or the issuer has to share the keys with the verifier). I don't really have a good argument for \emph{public} verifiability, however, given the state of VOPRFs, we can argue that we need the public verifiability. We need private verifiability e.g. in the case of PrivacyPass, where we need to ensure the issuer does not segment credentials by using different keys for a subset of parties. Right now, some deployments of PrivacyPass use blind RSA signatures, or ARC with BBS, but both are a lot more expensive than using a VOPRF. 
    \item[Encoding additional attributes] The credential structure in the presentation context should be flexible enough to allow encoding additional attributes beyond rate-limiting (e.g. origin, timestamps, validity period, metadata...) without requiring re-issuance. Especially \emph{late-origin binding} 
    \item[Binding the presentation limit] It's an open question whether the presentation limit $n$ should be cryptographically bound to the credential during issuance. Binding $n$ into the credential would prevent limit equivocation but reduce flexibility. If this is either hard or easy, it would be neat to form an opinion. 

\end{description}

\subsection{The generalized protocol}
\label{generalized_protocol}
To fix our construction in the right way, let's consider the following modified issuance and proof protocol. The message which will be signed during issuance will consist of two parts $\vectorify{m}_1$ and $\vectorify{m}_2$. $\vectorify{m}_1$ is used as actual public credential material that the user can reveal and prove something about \lena{What is the domain of $\vectorify{m}_1$?fl}. $\vectorify{m}_2\in \mathbb{F}_q^{\lenkey}$ is a uniformly random sequence  that we will later use as a key. In addition, we assume that the user uses a nonce when showing a credential. This nonce should be a value from  $D\subseteq \mathbb{F}_q^{\lennonce}$ for a public parameter $\lennonce$. It could just be a small number or some larger value, if we want the nonce to also contain some session context. \lena{Why not bind the session context to the public credential material?}

Towards constructing the credential scheme, we assume that there exists a secure pseudorandom function \[
\prf: \mathbb{F}_q^{\lenkey} \times \mathbb{F}_q^{\lennonce} \rightarrow \mathbb{F}_q^{\lentag}
\] 
such that $\lentag \geq \lceil \lambda/\log_2(q)\rceil$. We assume that $\prf$ can be efficiently computed in our ZK proof system. As implementation for $\prf$ we will likely use some conservative parameterization of Poseidon/Hades \cite{USENIX:GKRRS21,AFRICACRYPT:GraKhoSch23}. While usually running over larger fields, this should still work nicely in our setting. Though Farfalle~\cite{ToSC:BDHPVV17} might appear to be a more obvious choice, it unfortunately only works modulo $2$.
We now describe $\prf$ in more detail: 
Let $d\geq 3$ be the smallest constant such that $gcd(d,q-1)=1$. Let $R_F, R_P\in \mathbb{N}$ be fixed parameters. We define 
\[
\prf(\vectorify{k},\vectorify{n}):=Tr\left(P(\vectorify{k}\|\vectorify{n})+(\vectorify{k}\| \vectorify{n})\right)\text{.}
\]

Let $t=\lenkey+\lennonce$.
This definition uses the permutation $P: \mathbb{F}_q^{t} \rightarrow \mathbb{F}_q^{t}$ and the truncation function $Tr:\mathbb{F}_q^{t}\rightarrow \mathbb{F}_q^{\lentag}$ that simply outputs the first $\lentag$ elements. 

Now, we define the permutation $P$ which is basically the Poseidon2 permutation as defined in \cite{AFRICACRYPT:GraKhoSch23}.
It is defined as 
\[
P(\vectorify{x}):=\mathcal{E}_{R_F}\circ \cdots \circ \mathcal{E}_{R_F/2+1}\circ \mathcal{I}_{R_P}\circ \cdots \circ \mathcal{I}_{R_1}\circ \mathcal{E}_{R_F/2}\circ \cdots \circ \mathcal{E}_{R_1}(\vectorify{x})
\]
where $\mathcal{E}_{i}$ is the $i$th \emph{external round} of the permutation and $\mathcal{I}_{i}$ is the $i$th \emph{internal round}. \lena{explain external/internal round?}

For each $i\in[R_F],j\in [t]$ we fix round constants $c_j^{(i)}$. In addition, let $\hat{c}^{(i)}$ be fixed round constants for $i\in [R_P]$ and let $\vectorify{M}_{\mathcal{E}},\vectorify{M}_{\mathcal{I}}$ be MDS matrices.

Note that $\vectorify{M}_{\mathcal{E}},\vectorify{M}_{\mathcal{I}}$ are chosen ahead of time where we will set $\vectorify{M}[i,j]=\frac{1}{x_i+y_j}$  for mutually distinct $\{x_i\}_{i\in [t]},\{y_i\}_{i\in [t]}$ such that $x_i+y_j\neq 0$. This makes every submatrix of $\vectorify{M}$ invertible. In \cite{AFRICACRYPT:GraKhoSch23} the authors additionally require that $\vectorify{M}_{\mathcal{I}}$ must be chosen such that no arbitrarily long subspace trails exist. This holds if the minimal polynomials of $\vectorify{M}_{\mathcal{I}},\dots,\vectorify{M}_{\mathcal{I}}^{R_P}$ are irreducible and of maximum degree. Towards this, let a polynomial $P\in \mathbb{F}_q[X]$ be annihilating for a matrix $\vectorify{M}\in \mathbb{F}_q^{t\times t}$ if $P(\vectorify{M})\times \vectorify{v}=\vectorify{0}$ for all $\vectorify{v}\in \mathbb{F}_q^t$.
We define the minimal polynomial $\Phi\in \mathbb{F}_q[X]$ of a matrix $\vectorify{M}$ as the monic polynomial of minimal degree such that
\begin{enumerate}
    \item $\Phi$ is annihilating.
    \item For each polynomial $P\in \mathbb{F}_q[X]$ that is also annihilating, $\Phi$ divides $P$.
\end{enumerate}
\lena{why does it have to be annihilating?}
For security reasons, we need that $t-\lentag > \lambda/\log_2(q)$ as well to avoid a guessing attack on the output state. To achieve $\lambda$ bits of security, we must have $\lambda\approx  t\cdot \log_2(p)/3$. For now, we choose $R_F=8$ and $R_P=10$ as example parameters, and let $\lentag=2$, $\lennonce=8$ and $\lenkey=90$ \url{https://extgit.isec.tugraz.at/krypto/hadeshash/-/blob/master/code/generate_params_poseidon.sage} should give us code how to test if matrices fulfill these criteria. The algorithms are defined in~\cite{ToSC:GraRecSch21}. Note that choosing the matrices and round constants is done ahead of time and they are simply public system parameters.



We then define $$\mathcal{E}_{i}(\vectorify{x}):=\vectorify{M}_\mathcal{E}\times((\vectorify{x}[1]+c_1^{(i)})^d,\dots,(\vectorify{x}[t]+c_t^{(i)})^d)^\top
$$
and
$$\mathcal{I}_{i}(\vectorify{x}):=\vectorify{M}_\mathcal{I}\times((\vectorify{x}[1]+\hat{c}^{(i)})^d,\vectorify{x}[2],\dots,\vectorify{x}[t])^\top
$$


\paragraph{The actual credential scheme}
In the following, we will also use a function 
$$\centerfunc: [-2B,2B] \rightarrow [-B,B] \text{.}$$ Basically we want that for any two values $a,b\in [-B,B]$, if either is uniformly random then $\centerfunc(a+b)$ is uniformly random in $[-B,B]$. 
This can be achieved by ``wrapping around'' the interval boundaries back into the interval:
$$
\centerfunc(x)=\begin{cases}
    x & \text{ if } x\in [-B,B]\\
   x+2B+1 & \text { if } x\in [-2B,-B-1] \\
x-2B-1 & \text { if } x\in [B+1,2B] 
\end{cases}
$$

\begin{lemma}
Let $a\in [-B,B]$ be fixed. Then 
\[
\Pr_{b\in [-B,B],c=[-B,B]}\left[ c=\centerfunc(a+b) \right]=1/(2B+1)
\]
\end{lemma}
\begin{proof}
Should be trivial. Fix any $a \in [-B,B]$, then for any choice of $b\in [-B,B]$ we get the values $a-B,a-B+1,\dots,a+B-1,a+B$. For $a=0$ the statement clearly holds. If $a>0$ then the values $B+1,\dots,B+a$ are mapped by $\centerfunc$ to $-B,-B+a-1$ uniquely. However, $a-B$ is not changed by $\centerfunc$, so we map to the exact same interval. The same should go for $a<0$ due to symmetry.
\end{proof}


\begin{description}
   \item[Setup] Generate a public matrix $\vectorify{A}$ with a trapdoor kept by the Issuer. Also generate a random matrix $\vectorify{A}^{\committag}\in\in \R_q^{\lencom\times (\lenmsg+1+\lenrandpar{0}+\lenrandcom)}$ which can be part of the CRS. Here, $\lenmsg$ is the message length of the message of the user input to $h$, $\lenrandpar{0},1$ are the lengths of the randomnesses $\vectorify{r}_{0}^U,\vectorify{r}_{1}^U$ of our inputs to the hash function $h$ and $\lenrandcom$ is the length of an additional blinding randomness used in the Ajtai commitment created by $\vectorify{A}^{\committag}$. $\lencom$ is the length of the generated Ajtai commitment.
    \item[Issuance]~ The user starts with a message $\vectorify{m}_1\in [-B,B]^{\lenmsgpar{1}N}$ as input. Later on, we will let our proof system prove arbitrary statements about $\vectorify{m}_1$ as partial reveals.
    \begin{enumerate}
        \item The user samples $\vectorify{m}_2\in [-B,B]^{(\lenmsg-\lenmsgpar{1})N}$ as well as $\vectorify{r}_{0}^U\in [-B,B]^{\lenrandpar{0}N},\vectorify{r}_{1}^U\in [-B,B]^N$ uniformly at random. In addition, it chooses $\overline{\vectorify{r}}\in [-B,B]^{ \lenrandcom \cdot N}$.
        \item Let $\vectorify{A}^{\committag}$ be a uniformly random matrix. Compute $$\vectorify{c}\gets \vectorify{A}^{\committag}\begin{bmatrix}
            \vectorify{m}_1\\\vectorify{m}_2\\\vectorify{r}_{0}^U \\ \vectorify{r}_{1}^U \\ \overline{\vectorify{r}}
        \end{bmatrix}$$
        \item The user sends $\vectorify{c}$ to the issuer, together with a ZK proof of plaintext knowledge $\pi_{\committag}$ that it knows a preimage of $\vectorify{A}^{\committag}$. {\color{blue} this proof can likely be skipped, as we have to show the same thing again below.}\lena{ Agreed.}

    \item Issuer verifies $c$ using $\pi_{\committag}$ and rejects if the proof rejects. Otherwise it chooses $\vectorify{r}_{0}^I\in [-B,B]^{\lenrandpar{0}N},\vectorify{r}_{1}^I\in [-B,B]^N$ and sends these to the user.
    \item Upon receiving $\vectorify{r}_{0}^I,\vectorify{r}_{1}^I$ from the Issuer, the User sets \begin{align*}`\vectorify{r}_0&:=\centerfunc(\vectorify{r}_{0}^U+\vectorify{r}_{0}^I)\\
    \vectorify{r}_1&:=\centerfunc(\vectorify{r}_{1}^U+\vectorify{r}_{1}^I)\text{.}
    \end{align*}
        It then sets $\vectorify{m}:=\vectorify{m}_1\|\vectorify{m}_2$, computes $$ \vectorify{t}=h_{\vectorify{m},(\vectorify{r}_0,\vectorify{r}_1)}(\vectorify{B}) $$
        and sends $\vectorify{t}$ together with a proof $\pi_{\tagtag}$ to the Issuer. The proof shows that User knows $(\vectorify{m}_1,\vectorify{m}_2, \vectorify{r}_{0}^U,\vectorify{r}_{1}^U,\overline{\vectorify{r}})$  such that 
        \begin{enumerate}
            \item  $\vectorify{c} = \vectorify{A}^{\committag}(\vectorify{m}_1,\vectorify{m}_2, \vectorify{r}_{0}^U,\vectorify{r}_{1}^U,\overline{\vectorify{r}})^\top$
            \item $\vectorify{t}=h_{\vectorify{m}_1\|\vectorify{m}_2,\centerfunc(\vectorify{r}_{0}^U+\vectorify{r}_{0}^I),\centerfunc(\vectorify{r}_{1}^U+\vectorify{r}_{1}^I)}(\vectorify{B})$
            \item $\vectorify{m}_1,\vectorify{m}_2, \vectorify{r}_{0}^U,\vectorify{r}_{1}^U,\overline{\vectorify{r}}$ are from $[-B,B]$
        \end{enumerate}
  where $\vectorify{r}_{0}^I,\vectorify{r}_{1}^I$ are public.
        \item Upon verifying $\pi_{\tagtag}$, the issuer samples a preimage $\vectorify{u}$ under $\vectorify{A}$ and sends it to the User.
        \item The user checks if $\vectorify{t}=\vectorify{A}\vectorify{u}$ and that $\vectorify{u}$ is short. If so, then it accepts $\vectorify{u}$ as a valid signature.
    \end{enumerate}


    \item[Showing] To create a showing, the user chooses a previously unknown nonce $\acnonce\in D$ and computes the following:
    \begin{enumerate}
        \item Set $\actag=\prf(\vectorify{m}_2,\acnonce)$.
        \item Compute a Non-Interactive ZKPoK $\pi_{\showtag}$ showing knowledge of $\vectorify{u},\vectorify{m}_1,\vectorify{m}_2,\vectorify{r}_0,\vectorify{r}_1$ such that
        \begin{itemize}
            \item $\vectorify{A}\vectorify{u}= h_{\vectorify{m}_1\|\vectorify{m}_2,\vectorify{r}_0,\vectorify{r}_1}(\vectorify{B})$ 
            \item $\actag=\prf(\vectorify{m}_2,\acnonce)$
            \item $\vectorify{m}_1,\vectorify{m}_2,\vectorify{r}_0,\vectorify{r}_1$ are over $[-B,B]$
        \end{itemize}
        \item Then, send $\actag,\acnonce,\pi_{\showtag}$ to the verifier. \lena{do we want to add a range proof variant? }
        \item The verifier accepts the proof $\pi_{\showtag}$ if it verifies and if $\acnonce,\actag$ has not been revealed before. {\color{blue} could likely also just show that nonce is from some interval and keep it secret as well.}
    \end{enumerate}
\end{description}


\subsection{ARC Implementation}


\paragraph{Notation and naming conventions.}
In the implementation we distinguish the public hash key $B$ used by $h(\cdot;B)$ from the coefficient bound used
for sampling and membership constraints. Concretely:
\begin{itemize}
  \item \texttt{$\BoundMsg$} is the coefficient bound used in all intervals and in the centering map.
  \item $B$ (and its cleared-denominator pieces $B_0,B_1,B_2,B_3$) is the public vSIS/BBS hash key.
\end{itemize}
So whenever the generalized description writes an interval like $[-B,B]$ as a shortness bound, the implementation
uses $[-\BoundMsg,\BoundMsg]$.

\subsubsection{The commitment scheme and helper functions}
We implement the Ajtai-style linear commitment. The commitment is
\[
\mathrm{com}=A_c\cdot[m_1\;\|\;m_2\;\|\;r^U_0\;\|\;r^U_1\;\|\;\bar r]^\top,
\]
where $A_c$ is a public random matrix and all operands are ring polynomials.
All matrices and vectors live in the same ring/field as the rest of the protocol: the cyclotomic ring
\[
R_q=\mathbb{Z}_q[X]/(X^N+1)\,.
\]
Concretely, $A_c$, $B$, and the signature matrix $A$ are represented as matrices of ring polynomials over $\mathbb{F}_q$
(coefficients in $[0,q)$), and all vector entries $m_1,m_2,r^U_0,r^U_1,\bar r,r_0,r_1$ are ring elements in the same
domain.

\lena{What is the domain of $\vectorify{m}_1$?fl}
% Answer: In the concrete implementation, m1 is a bounded ring vector. We sample it coefficient-wise in
% [-BoundB,BoundB], and the PACS statement enforces the same membership constraint.

We implement the \texttt{center} map via \texttt{CenterBounded}, which wraps sums back into
$[-\BoundMsg,\BoundMsg]$ using $\Delta=2\cdot\BoundMsg+1$. The protocol uses this to combine Holder
and Issuer randomness without bias:
\[
r_0=\mathrm{center}(r^U_0+r^I_0),\qquad r_1=\mathrm{center}(r^U_1+r^I_1).
\]
In code, \texttt{CombineRandomness} applies \texttt{CenterBounded} coefficient-wise and records carry terms (K0/K1) used
by the center constraints in the NIZK.
\subsubsection{The PRF}
We implement the PRF from \ref{generalized_protocol} exactly as specified: a Poseidon2-like permutation with
feed-forward and truncation. The PRF is
\[
\mathrm{tag}=\mathrm{Tr}\big(P(k\|n)+ (k\|n)\big),
\]
where $k=m_2$ is the secret key vector, $n$ is the public nonce, and $t=\lenkey+\lennonce$ is the permutation width.
The truncation $\mathrm{Tr}$ outputs the first $\lentag$ lanes.

In the codebase, $m_2$ is represented as bounded ring elements (so it fits the same boundedness constraints as the rest
of the credential witness). When we need a vector in $\mathbb{F}_q^{\lenkey}$, we use a fixed public encoding that
extracts the relevant lanes; the showing proof binds those key lanes back to the signed $m_2$.

\lena{Why not bind the session context to the public credential material?}
% Answer: If we want session binding, we do it at showing time by extending the PRF input (nonce || ctx).
% We avoid embedding session context into the signed public attributes, since that would make the credential
% session-specific and hurt reusability.

\paragraph{Parameters and generation.}
All PRF parameters are public and loaded from \texttt{prf/prf\_params.json}. This file is generated by the Sage script
\texttt{prf/generate\_params\_poseidon.sage}, which fixes:
\begin{itemize}
  \item field modulus $q$ (same $q$ as the PCS ring),
  \item S-box exponent $d$ (smallest $d\ge 3$ with $\gcd(d,q-1)=1$),
  \item width $t=\lenkey+\lennonce$ and truncation length $\lentag$,
  \item round counts $R_F$ (external) and $R_P$ (internal),
  \item MDS matrices $M_E,M_I$ (Cauchy-style, checked for invertibility),
  \item round constants $c^{(i)}$ (vector per external round) and $\hat c^{(i)}$ (scalar per internal round).
\end{itemize}
The loader in \texttt{prf/params.go} validates dimensions and the evenness of $R_F$. The Sage script also reports the
heuristic permutation security and the truncation bound $t-\lentag>\lambda/\log_2(q)$ so parameters can be tuned.

\paragraph{Permutation construction.}
The permutation $P$ is the composition
\[
P=\mathcal{E}_{R_F}\circ \cdots \circ \mathcal{E}_{R_F/2+1}\circ \mathcal{I}_{R_P}\circ \cdots \circ \mathcal{I}_1\circ
\mathcal{E}_{R_F/2}\circ \cdots \circ \mathcal{E}_1.
\]
\lena{explain external/internal round?}
% Answer: External rounds are full rounds (S-box on all lanes). Internal rounds are partial rounds
% (S-box only on lane 1). This reduces nonlinear constraints in the ZK arithmetization while keeping diffusion.

Each external round $\mathcal{E}_i$ adds the per-lane constants, applies the S-box to every lane, and multiplies by $M_E$:
\[
\mathcal{E}_i(x)=M_E\cdot\big((x_1+c^{(i)}_1)^d,\dots,(x_t+c^{(i)}_t)^d\big)^\top.
\]
Each internal round $\mathcal{I}_i$ adds a scalar constant to lane 1, applies the S-box only to lane 1, keeps the other
lanes linear, and multiplies by $M_I$:
\[
\mathcal{I}_i(x)=M_I\cdot\big((x_1+\hat c^{(i)})^d,x_2,\dots,x_t\big)^\top.
\]
We implement these steps in \texttt{prf/permute.go} using the same modulus $q$ as the rest of the protocol.

\paragraph{Implementation details.}
The Go API \texttt{prf.Tag} builds the state $x=k\|n$, applies \texttt{PermuteInPlace}, performs the feed-forward
$x\leftarrow x+ x^{(0)}$, and returns the first $\lentag$ lanes.

For proving, we use a \emph{nonlinear-only} PRF witness: we commit only the PRF key lanes and selected S-box outputs,
and we recompute all linear wiring (add-round-constants and MDS multiplications) inside the constraint evaluator.
Concretely, \texttt{prf/trace.go} provides:
\begin{itemize}
  \item \texttt{TraceSBoxOutputs} and \texttt{TraceSBoxOutputsGrouped}, which return only the (checkpointed) S-box outputs in
  strict execution order (used by the showing NIZK).
\end{itemize}
All arithmetic is performed in the base field $\mathbb{F}_q$ (same modulus as the PCS ring), so PRF outputs and
parameters are directly compatible with the constraint system.

\subsubsection{The Pre-Sign NIZK and Issuance Protocol}
The issuance (pre-sign) proof instantiates the statement in Section B.6 as a PACS/PIOP instance over the witness matrix
$w\in\mathbb{F}_q^{n\times s}$, where each row is a ring element encoded as a polynomial $P_i(X)$ with
\[
P_i(\omega_k)=w_{i,k}\quad(k\in[1,s]),
\]
and degree bounded by $s+\ell-1$ after adding $\ell$ random tail points. We use one row per ring element, with the
following fixed order (single-poly blocks):
\[
M1,\,M2,\,RU0,\,RU1,\,R,\,R0,\,R1,\,K0,\,K1.
\]
Here $M1,M2$ are the message halves, $RU0,RU1,R$ are Holder randomnesses, $R0,R1$ are the centered combinations, and
$K0,K1$ are the center carry rows (coefficients in $\{-1,0,1\}$). The target $T$ is \emph{public} in the pre-sign
statement and is therefore not a witness row.

\paragraph{Public inputs.}
The verifier receives $A_c$, $B$, $\mathrm{com}$, $r^I_0$, $r^I_1$, and $T$, all represented as public polynomials over
the same ring:
\begin{itemize}
  \item $\mathrm{com}$ is the commitment vector (rows of $A_c\cdot[M1\|M2\|RU0\|RU1\|R]^\top$).
  \item $r^I_0,r^I_1$ are issuer randomness polynomials.
  \item $B=(B_0,B_1,B_2,B_3)$ is the vSIS/BBS hash key matrix in the cleared-denominator form.
  \item $T$ is the public hash target (one polynomial).
\end{itemize}
All public parameters are interpolated into $\Theta$-polynomials over $\Omega$ so that constraints are evaluated in the
same domain as the witness rows.

\lena{I'm not sure where $B$ is coming from, and how we can prove efficiently that the values are the same as in the commitment. Has this challenge approach been used in lattices before? It seems quite straightforward in any case.}
% Answer: B is a public setup parameter (hash key), as in vSIS/BBS. Consistency is enforced by the single pre-sign PACS
% proof: we check the commitment opening and the hash/centering relation in the same statement, so T is tied to the
% committed witness. The issuer randomness (r^I_0,r^I_1) is the standard de-biasing step used in interactive variants.

\paragraph{Constraint families (all parallel, F\textsubscript{par}).}
\begin{enumerate}
  \item \textbf{Commit constraints (linear):}
  \[
  A_c\cdot[M1\|M2\|RU0\|RU1\|R]^\top = \mathrm{com}.
  \]
  For each commitment output row $\ell$ and each $\omega_k$, we evaluate a linear residual
  \[
  f^{\mathrm{com}}_\ell(\omega_k)=\sum_j (A_c)_{\ell,j}(\omega_k)\cdot w_{j,k}-\mathrm{com}_\ell(\omega_k).
  \]

  \item \textbf{Center constraints (linear with carries):}
  Using $\Delta=2\cdot\BoundMsg+1$,
  \[
  R0 = RU0 + r^I_0 - \Delta\cdot K0,\qquad R1 = RU1 + r^I_1 - \Delta\cdot K1,
  \]
  evaluated slot-wise. The carry rows $K0,K1$ are bounded to $\{-1,0,1\}$ via membership constraints (see below).

  \item \textbf{Hash constraints (cleared denominator, linear in pre-sign):}
  With public $T$, we enforce
  \[
  (B_3 - R1)\odot T - \big(B_0 + B_1\cdot(M1{+}M2) + B_2\cdot R0\big)=0,
  \]
  evaluated per coordinate on $\Omega$. Because $T$ is public in issuance, this relation is linear in the witness rows.

  \item \textbf{Packing constraints (message split):}
  $M1$ occupies the lower half of coefficients and $M2$ the upper half. We enforce
  \[
  M1(\omega_k)=0 \text{ for } k> s/2,\qquad M2(\omega_k)=0 \text{ for } k\le s/2.
  \]

  \item \textbf{Bounds (membership polynomial):}
  For each witness row $X\in\{M1,M2,RU0,RU1,R,R0,R1\}$ we enforce
  \[
  P_{\BoundMsg}(X(\omega_k))=0\quad\forall \omega_k\in\Omega,
  \]
  where $P_{\BoundMsg}$ is the vanishing polynomial over $[-\BoundMsg,\BoundMsg]$. For carry rows
  $K0,K1$, we use $P_1$ (vanishing over $\{-1,0,1\}$).
\end{enumerate}

\paragraph{Witness construction.}
The prover samples $M1,M2,RU0,RU1,R$ in coefficient form (bounded), derives $R0,R1$ via \texttt{CenterBounded}, and sets
$K0,K1$ to the wrap-around carries so that the center constraints hold exactly. The witness rows are then interpolated to
$\Omega$ with fresh $\ell$-length tail masks (LVCS) to obtain the committed witness polynomials $P_i(X)$. Because all
constraints are parallel, the verifier replay evaluates the above residuals at the Fiat-Shamir chosen points using the
opened row evaluations.

\paragraph{Current demo metrics (Jan 2026).}
On the current implementation (\texttt{go run ./cmd/issuance}), the pre-sign proof reports:
\begin{itemize}
  \item base rows $=9$ (no PRF rows), mask rows $=2$,
  \item proof size $\approx 7.43$\,KB ($7606$ bytes),
  \item prover time $\approx 0.028$\,s, verifier time $\approx 0.032$\,s,
  \item soundness bits $\approx 129.61$,
  \item parameter tuple $(\mathrm{NCols},\ell,\ell',\rho,\theta,\eta,d_Q)=(16,25,2,2,4,19,11575)$.
\end{itemize}

\subsubsection{The Post-Sign Showing Protocol}
The showing proof (post-sign) is a PACS/PIOP statement that proves, in zero-knowledge, that a holder possesses a valid
credential and that the public tag is correctly derived. It binds together: the signature relation, the hash relation,
the PRF relation, and coefficient-wise bounds. All constraints are evaluated on the same evaluation set $\Omega$ of size
$s=\mathrm{NCols}$, with witness polynomials of degree $\le s+\ell-1$ after LVCS masking. We operate in the small-field
path ($\theta>1$), so the verifier replays Eq.(4) using K-points and tail openings.

\paragraph{Witness rows and ordering.}
We commit a single witness matrix $w\in\mathbb{F}_q^{n\times s}$ whose rows are fixed-order ring elements. For showing we
use the following row layout:
\begin{enumerate}
  \item \textbf{Base credential rows} (same packing as issuance):
  \[
  M1,\;M2,\;RU0,\;RU1,\;R,\;R0,\;R1,\;K0,\;K1.
  \]
  These are single-poly blocks, with $M1$ and $M2$ packed into lower/upper halves of the ring (half-split).
  \item \textbf{Internal hash target:} $T$ (one row), treated as a witness in showing.
  \item \textbf{Signature preimage:} $U$ (vector; in our NTRU signature format this is 2 rows), where $A\cdot U = T$.
  \item \textbf{PRF witness rows (System A, grouped checkpoints):}
  after the base rows we append the PRF block:
  \[
  \mathrm{Key}[0],\dots,\mathrm{Key}[\lenkey-1],\qquad Z[0],\dots,Z[ZCount-1],
  \]
  where $\mathrm{Key}[i]$ are the PRF key lanes (one row per lane) and $Z[\alpha]$ are the \emph{checkpointed} S-box outputs
  in strict execution order. The PRF start index is recorded so constraints bind to the correct rows.
\end{enumerate}
Let $n_{\mathrm{base}}$ denote the number of base rows (including $T$ and $U$), and let $ZCount$ denote the number of
committed S-box outputs under the checkpoint schedule. Then
\[
n_{\mathrm{prf}}=\lenkey + ZCount,\qquad n=n_{\mathrm{base}}+n_{\mathrm{prf}}.
\]
All PRF witness rows are constant over $\Omega$ (column-constant), but still masked by the LVCS tail to preserve ZK.

\paragraph{Concrete indices (current instantiation).}
With the NTRU signature format used in this repo, $U$ has two rows $(s_0,s_1)$, so
\[
n_{\mathrm{base}} = 9\;(\text{base}) + 1\;(\text{T}) + 2\;(\text{U}) = 12.
\]
The row indices are therefore:
\[
0..8:\;M1,M2,RU0,RU1,R,R0,R1,K0,K1;\quad 9:\;T;\quad 10..11:\;U.
\]
The PRF block begins at $\texttt{startIdx}=12$ and has $n_{\mathrm{prf}}=\lenkey+ZCount$ rows. For the current default
PRF parameters $(t,R_F,R_P)=(20,8,18)$ and grouping $g=2$ (checkpoint every 2nd full round except the last; internal rounds
always checkpointed), we have:
\[
ZCount = R_P + 3\cdot t = 18 + 60 = 78,\qquad n_{\mathrm{prf}}=\lenkey+ZCount=8+78=86,
\]
so the total witness rows are $n=12+86=98$. These counts are exposed in logs during \texttt{cmd/showing} to confirm the
split between base and PRF rows.

\paragraph{Public inputs.}
The verifier receives:
\begin{itemize}
  \item $A$ (signature matrix), $B$ (hash key matrix),
  \item $\mathrm{tag}$ and $\mathrm{nonce}$ (both public),
  \item $\BoundMsg$ (coefficient bound),
  \item optionally $\mathrm{Com}$ and $A_c$ if re-binding to issuance is desired.
\end{itemize}
Public matrices and vectors are interpolated into $\Theta$-polynomials over $\Omega$ so constraints are evaluated in the
same domain as the witness rows.

\paragraph{Constraint families (F\textsubscript{par} only).}
All showing constraints are parallel; the aggregated set remains empty. Each residual is enforced slot-wise on $\Omega$
and recomputed by the verifier at Fiat-Shamir chosen points using opened row evaluations.

\begin{enumerate}
  \item \textbf{Signature constraint (linear):}
  \[
  A\cdot U = T.
  \]
  For each output row of $A$ and each $\omega_k$, we evaluate a linear residual
  \[
  f^{\mathrm{sig}}(\omega_k)=\sum_j A_{j}(\omega_k)\cdot U(\omega_k)-T(\omega_k).
  \]

  \item \textbf{Hash constraint (cleared denominator, bilinear):}
  With $T$ now a witness, the cleared-denominator hash becomes bilinear:
  \[
  (B_3 - R1)\odot T - \big(B_0 + B_1\cdot(M1{+}M2) + B_2\cdot R0\big)=0.
  \]

  \item \textbf{PRF constraints (System A, effective degree $d^g$):}
  The PRF constraints are expressed entirely as parallel constraints and are recomputed by the verifier from opened row
  evaluations and public $\Theta$ values:
  \begin{itemize}
    \item \emph{Key binding:} enforce that $\mathrm{Key}[i]$ equals the signed $m_2$ encoding carried by the $M2$ row at the
    designated packing coordinate on $\Omega$ (half-split packing). We implement this using a selector $\Theta$ polynomial
    that is $1$ at that coordinate and $0$ elsewhere, so the equality is enforced at exactly one $\Omega$ slot per key lane.
    \item \emph{Checkpointed S-box constraints:} for each committed $Z_\alpha$, enforce $Z_\alpha - U_\alpha^d=0$ where
    $U_\alpha$ is the computed S-box input lane value at that checkpoint (linear wiring is recomputed inside the evaluator).
    \item \emph{Tag binding (feed-forward + truncation):} for each lane $j<\lentag$, enforce $\mathrm{tag}_j = y_j + x^{(0)}_j$,
    where $y=P(x^{(0)})$ is the final state.
  \end{itemize}
  With grouping $g=2$ we effectively compose one uncheckpointed nonlinear layer, so the resulting constraints have
  effective degree $\approx d^g=25$ in the underlying witness variables, while still being evaluated by the verifier as
  public residual functions on opened row evaluations (paper-faithful Eq.(4) replay).

  \item \textbf{Packing constraints:}
  \[
  M1(\omega_k)=0 \text{ for } k> s/2,\qquad M2(\omega_k)=0 \text{ for } k\le s/2.
  \]

  \item \textbf{Bounds:}
  For each bounded row $X\in\{M1,M2,R0,R1,T,U\}$ (and other base rows if needed),
  \[
  P_{\BoundMsg}(X(\omega_k))=0\quad\forall \omega_k\in\Omega,
  \]
  with $P_{\BoundMsg}$ the vanishing polynomial on $[-\BoundMsg,\BoundMsg]$. For carry rows $K0,K1$, we use
  $P_1$ (vanishing on $\{-1,0,1\}$).
\end{enumerate}

\paragraph{Column count, degrees, and $d_Q$.}
The evaluation domain has size $s=\mathrm{NCols}$. Each witness row is interpolated with $\ell$ random tail points, so
\[
\deg(P_i)\le s+\ell-1.
\]
The effective maximum constraint degree for showing is dominated by:
\begin{itemize}
  \item the PRF S-box exponent composed across checkpoint groups ($d^g$, with $d=5$ and $g=2$ giving $\approx 25$), and
  \item the membership polynomial degree $\deg(P_{\BoundMsg})=2\cdot\BoundMsg+1$ used for coefficient bounds
  (for \BoundMsg=8, $\deg(P_{\BoundMsg})=17$).
\end{itemize}
In the implementation, the masking degree bound $d_Q$ is computed by \texttt{deriveMaskingConfig} from the instantiated
ring-polynomial degrees of the constructed $F$-polynomials (mirroring the PACS path) and is reported in the showing logs.

\paragraph{Verifier replay (small-field, $\theta>1$).}
For $\theta>1$, the verifier replays Eq.(4) using K-point openings. It receives masked evaluations of the committed rows
at Fiat-Shamir chosen K-points, computes the parallel residuals from those row evaluations and public $\Theta$ values,
and checks that the reconstructed $Q_i(e)$ matches the prover's masked combination (and that
$\sum_{\omega\in\Omega}Q_i(\omega)=0$). This binds the showing proof to the committed witness rows while preserving
zero-knowledge via LVCS/DECS masking.

\paragraph{Current demo metrics (Jan 2026).}
On the current implementation (\texttt{go run ./cmd/showing -prfGroupRounds 2}), the showing proof reports:
\begin{itemize}
  \item base rows $=12$, PRF rows $=86$ (total $98$), mask rows $=2$,
  \item proof size $\approx 19.96$\,KB ($20443$ bytes),
  \item prover time $\approx 0.323$\,s, verifier time $\approx 0.053$\,s,
  \item soundness bits $\approx 129.44$,
  \item parameter tuple $(\mathrm{NCols},\ell,\ell',\rho,\theta,\eta,d_Q)=(16,25,2,2,4,19,17015)$.
\end{itemize}
